
\documentclass[11pt,a4paper]{article}
\usepackage[utf8]{inputenc}
\usepackage[T1]{fontenc}
\usepackage[french]{babel}
\usepackage{lmodern}
\usepackage{microtype}
\usepackage[a4paper,margin=1.65cm,headheight=15pt]{geometry}
\usepackage{amsmath,amssymb,mathtools}
\usepackage{array,tabularx,multicol,booktabs}
\usepackage{enumitem}
\usepackage[most]{tcolorbox}
\usepackage{xcolor}
\usepackage{tikz}
\usetikzlibrary{arrows.meta,calc,positioning}
\usepackage{pdflscape}
\usepackage{fancyhdr}
\usepackage{hyperref}
\usepackage{needspace}
\usepackage{pagecolor}

\definecolor{fondpage}{RGB}{247,248,250}
\definecolor{encre}{RGB}{31,35,40}
\definecolor{bleu}{RGB}{40,91,135}
\definecolor{vert}{RGB}{55,125,92}
\definecolor{orange}{RGB}{178,105,42}
\definecolor{violet}{RGB}{101,77,145}
\definecolor{rouge}{RGB}{170,72,72}
\definecolor{gris}{RGB}{86,91,99}
\definecolor{bleufonce}{RGB}{28,55,120}
\definecolor{bleuclair}{RGB}{238,243,255}
\definecolor{grisclair}{RGB}{245,245,245}
\definecolor{tableline}{RGB}{45,45,45}
\definecolor{softgray}{RGB}{246,247,249}
\definecolor{deepblue}{RGB}{25,62,115}
\definecolor{accent}{RGB}{20,82,130}

\hypersetup{colorlinks=true,linkcolor=bleufonce,urlcolor=bleufonce,
pdftitle={Parcours de revision -- Dérivation -- Premiere specialite},pdfauthor={}}

\setlength{\parindent}{0pt}
\setlength{\parskip}{0.25em}
\renewcommand{\arraystretch}{1.13}
\setlength{\tabcolsep}{4pt}
\setlist[itemize]{leftmargin=1.25em,itemsep=0.15em,topsep=0.2em}
\setlist[enumerate,1]{label=\textbf{\arabic*.}, leftmargin=1.05cm, itemsep=0.35em, topsep=0.2em}
\setlist[enumerate,2]{label=\textbf{\alph*.}, leftmargin=0.85cm, itemsep=0.25em, topsep=0.15em}

\pagestyle{fancy}
\fancyhf{}
\cfoot{\thepage}

\newcommand{\R}{\mathbb{R}}
\newcommand{\N}{\mathbb{N}}
\newcommand{\sourceq}[1]{ {\scriptsize\textcolor{bleufonce}{(site Première q#1)}}}
\newcommand{\code}[1]{\texttt{#1}}

\newcounter{question}
\newcounter{reponse}

\newenvironment{blocquestions}[1]{%
  \begin{tcolorbox}[enhanced,breakable,colback=bleuclair,colframe=bleufonce,boxrule=0.6pt,arc=2mm,left=2mm,right=2mm,top=2mm,bottom=1.5mm,title={\bfseries #1},coltitle=white,colbacktitle=bleufonce,before skip=3mm,after skip=3mm, fontupper=\large]
  \large
  \begin{enumerate}[leftmargin=*,label=\textbf{\arabic*.},itemsep=0.82em,topsep=0.5em]
  \setcounter{enumi}{\value{question}}
}{%
  \setcounter{question}{\value{enumi}}
  \end{enumerate}
  \end{tcolorbox}
}

\newenvironment{blocreponses}[1]{%
  \begin{tcolorbox}[enhanced,breakable,colback=bleuclair,colframe=bleufonce,boxrule=0.6pt,arc=2mm,left=2mm,right=2mm,top=2mm,bottom=1.5mm,title={\bfseries #1},coltitle=white,colbacktitle=bleufonce,before skip=3mm,after skip=3mm, fontupper=\large]
  \large
  \begin{enumerate}[leftmargin=*,label=\textbf{\arabic*.},itemsep=0.42em,topsep=0.5em]
  \setcounter{enumi}{\value{reponse}}
}{%
  \setcounter{reponse}{\value{enumi}}
  \end{enumerate}
  \end{tcolorbox}
}

\newtcolorbox{correctionbox}[1]{enhanced,breakable,colback=grisclair,colframe=black!55,boxrule=0.55pt,arc=2mm,left=2mm,right=2mm,top=2mm,bottom=1.5mm,title=\textbf{#1},coltitle=black,colbacktitle=black!12,before skip=2mm,after skip=2mm}

\newtcolorbox{methodebox}[1]{enhanced,breakable,colback=white,colframe=bleu!70!black,boxrule=0.55pt,arc=2mm,left=2mm,right=2mm,top=2mm,bottom=1.5mm,title=\textbf{#1},coltitle=white,colbacktitle=bleu!70!black,before skip=2mm,after skip=2mm}

\newcommand{\titreSujet}[2]{%
  \subsection*{#1}
  \addcontentsline{toc}{subsection}{#1}
  \textit{Référence / inspiration : #2}\par\medskip\hrule\medskip
}

\newcommand{\qcm}[4]{%
\begin{center}
\begin{tabularx}{0.96\linewidth}{@{}XX@{}}
\textbf{a.} #1 & \textbf{b.} #2\\[1mm]
\textbf{c.} #3 & \textbf{d.} #4
\end{tabularx}
\end{center}}

\tcbset{enhanced,arc=2mm,outer arc=2mm,boxrule=0.42pt,left=1.15mm,right=1.15mm,top=0.75mm,bottom=0.75mm,boxsep=0.35mm,before skip=0.8mm,after skip=0.8mm,fonttitle=\bfseries\footnotesize,coltitle=encre,before upper=\raggedright\fontsize{7.65}{8.15}\selectfont}
\newtcolorbox{fichebox}[3][]{colback=white,colframe=#2!72!black,colbacktitle=#2!18,coltitle=#2!50!black,borderline west={0.9mm}{0pt}{#2!78!black},title={#3},drop fuzzy shadow=black!5,#1}
\newcommand{\tagcol}[2]{\begin{tcolorbox}[enhanced,colback=#1!11,colframe=#1!45!black,boxrule=0.42pt,arc=2mm,left=1mm,right=1mm,top=0.45mm,bottom=0.45mm,before skip=0mm,after skip=0.85mm,fontupper=\bfseries\scriptsize,colupper=#1!55!black]\centering #2\end{tcolorbox}}
\newtcolorbox{panel}[1]{colback=fondpage,colframe=fondpage,boxrule=0pt,arc=0mm,left=0mm,right=0mm,top=0mm,bottom=0mm,height=168mm,valign=top,before skip=0mm,after skip=0mm}
\newcommand{\titrepage}{\begin{tcolorbox}[colback=white,colframe=bleu!70!black,boxrule=0.65pt,arc=3mm,left=2.5mm,right=2.5mm,top=1.15mm,bottom=1.15mm,before skip=0mm,after skip=1.4mm,drop fuzzy shadow=black!10]
{\Large\bfseries Parcours de révision -- Dérivation}\hfill {\normalsize Première spécialité -- sans calculatrice}
\end{tcolorbox}}

% Tableaux TikZ sobres pour les corrections
\newcommand{\TableauSigneDerivDeux}{%
\begin{center}
\begin{tikzpicture}[x=1cm,y=0.62cm,>=Latex]
\def\L{1.8}\def\W{12.4}\def\H{2.35}
\fill[black!8] (0,1.55) rectangle (\W,\H);
\fill[black!5] (0,0) rectangle (\L,\H);
\draw[tableline, line width=0.45pt] (0,0) rectangle (\W,\H);
\draw[tableline, line width=0.45pt] (\L,0)--(\L,\H);
\draw[tableline, line width=0.45pt] (0,1.55)--(\W,1.55);
\draw[tableline, line width=0.35pt] (5.2,0)--(5.2,\H);
\draw[tableline, line width=0.35pt] (8.6,0)--(8.6,\H);
\node at (0.9,1.95) {$x$};
\node at (0.9,0.75) {$f'(x)$};
\node at (2.25,1.95) {$-\infty$};
\node at (5.2,1.95) {$-3$};
\node at (8.6,1.95) {$1$};
\node at (12.0,1.95) {$+\infty$};
\node at (3.55,0.75) {$+$};
\node at (5.2,0.75) {$0$};
\node at (6.9,0.75) {$-$};
\node at (8.6,0.75) {$0$};
\node at (10.35,0.75) {$+$};
\end{tikzpicture}
\end{center}}

\newcommand{\TableauVariationDeux}{%
\begin{center}
\begin{tikzpicture}[x=1cm,y=0.62cm,>=Latex]
\def\L{1.8}\def\W{12.4}\def\H{3.65}
\fill[black!8] (0,2.82) rectangle (\W,\H);
\fill[black!5] (0,0) rectangle (\L,\H);
\draw[tableline, line width=0.45pt] (0,0) rectangle (\W,\H);
\draw[tableline, line width=0.45pt] (\L,0)--(\L,\H);
\draw[tableline, line width=0.45pt] (0,2.82)--(\W,2.82);
\draw[tableline, line width=0.45pt] (0,1.82)--(\W,1.82);
\draw[tableline, line width=0.35pt] (5.2,0)--(5.2,\H);
\draw[tableline, line width=0.35pt] (8.6,0)--(8.6,\H);
\node at (0.9,3.22) {$x$};
\node at (0.9,2.32) {$f'(x)$};
\node at (0.9,0.9) {$f$};
\node at (2.25,3.22) {$-\infty$};
\node at (5.2,3.22) {$-3$};
\node at (8.6,3.22) {$1$};
\node at (12.0,3.22) {$+\infty$};
\node at (3.55,2.32) {$+$};\node at (5.2,2.32) {$0$};\node at (6.9,2.32) {$-$};\node at (8.6,2.32) {$0$};\node at (10.35,2.32) {$+$};
\node at (2.25,0.35) {};\node at (5.2,1.45) {$29$};\node at (8.6,0.35) {$-3$};\node at (12.0,1.45) {};
\draw[->,line width=0.65pt] (2.35,0.35)--(4.8,1.35);
\draw[->,line width=0.65pt] (5.6,1.35)--(8.2,0.45);
\draw[->,line width=0.65pt] (9.0,0.45)--(11.55,1.35);
\end{tikzpicture}
\end{center}}

\newcommand{\TableauSigneAire}{%
\begin{center}
\begin{tikzpicture}[x=1cm,y=0.62cm,>=Latex]
\def\L{1.8}\def\W{10.8}\def\H{2.35}
\fill[black!8] (0,1.55) rectangle (\W,\H);
\fill[black!5] (0,0) rectangle (\L,\H);
\draw[tableline, line width=0.45pt] (0,0) rectangle (\W,\H);
\draw[tableline, line width=0.45pt] (\L,0)--(\L,\H);
\draw[tableline, line width=0.45pt] (0,1.55)--(\W,1.55);
\draw[tableline, line width=0.35pt] (6.3,0)--(6.3,\H);
\node at (0.9,1.95) {$x$};
\node at (0.9,0.75) {$A'(x)$};
\node at (2.25,1.95) {$0$};
\node at (6.3,1.95) {$3$};
\node at (10.35,1.95) {$6$};
\node at (4.25,0.75) {$+$};
\node at (6.3,0.75) {$0$};
\node at (8.25,0.75) {$-$};
\end{tikzpicture}
\end{center}}

\newcommand{\TableauVariationAire}{%
\begin{center}
\begin{tikzpicture}[x=1cm,y=0.62cm,>=Latex]
\def\L{1.8}\def\W{10.8}\def\H{3.65}
\fill[black!8] (0,2.82) rectangle (\W,\H);
\fill[black!5] (0,0) rectangle (\L,\H);
\draw[tableline, line width=0.45pt] (0,0) rectangle (\W,\H);
\draw[tableline, line width=0.45pt] (\L,0)--(\L,\H);
\draw[tableline, line width=0.45pt] (0,2.82)--(\W,2.82);
\draw[tableline, line width=0.45pt] (0,1.82)--(\W,1.82);
\draw[tableline, line width=0.35pt] (6.3,0)--(6.3,\H);
\node at (0.9,3.22) {$x$};
\node at (0.9,2.32) {$A'(x)$};
\node at (0.9,0.9) {$A$};
\node at (2.25,3.22) {$0$};
\node at (6.3,3.22) {$3$};
\node at (10.35,3.22) {$6$};
\node at (4.25,2.32) {$+$};\node at (6.3,2.32) {$0$};\node at (8.25,2.32) {$-$};
\node at (2.25,0.35) {$0$};\node at (6.3,1.45) {$18$};\node at (10.35,0.35) {$0$};
\draw[->,line width=0.65pt] (2.75,0.45)--(5.8,1.35);
\draw[->,line width=0.65pt] (6.8,1.35)--(9.9,0.45);
\end{tikzpicture}
\end{center}}


\newcommand{\TableauSigneExpo}{%
\begin{center}
\begin{tikzpicture}[x=1cm,y=0.62cm,>=Latex]
\def\L{1.8}\def\W{10.8}\def\H{2.35}
\fill[black!8] (0,1.55) rectangle (\W,\H);
\fill[black!5] (0,0) rectangle (\L,\H);
\draw[tableline, line width=0.45pt] (0,0) rectangle (\W,\H);
\draw[tableline, line width=0.45pt] (\L,0)--(\L,\H);
\draw[tableline, line width=0.45pt] (0,1.55)--(\W,1.55);
\draw[tableline, line width=0.35pt] (6.3,0)--(6.3,\H);
\node at (0.9,1.95) {$x$};
\node at (0.9,0.75) {$g'(x)$};
\node at (2.25,1.95) {$-2$};
\node at (6.3,1.95) {$0$};
\node at (10.35,1.95) {$3$};
\node at (4.25,0.75) {$-$};
\node at (6.3,0.75) {$0$};
\node at (8.25,0.75) {$+$};
\end{tikzpicture}
\end{center}}

\newcommand{\TableauVariationExpo}{%
\begin{center}
\begin{tikzpicture}[x=1cm,y=0.62cm,>=Latex]
\def\L{1.8}\def\W{10.8}\def\H{3.65}
\fill[black!8] (0,2.82) rectangle (\W,\H);
\fill[black!5] (0,0) rectangle (\L,\H);
\draw[tableline, line width=0.45pt] (0,0) rectangle (\W,\H);
\draw[tableline, line width=0.45pt] (\L,0)--(\L,\H);
\draw[tableline, line width=0.45pt] (0,2.82)--(\W,2.82);
\draw[tableline, line width=0.45pt] (0,1.82)--(\W,1.82);
\draw[tableline, line width=0.35pt] (6.3,0)--(6.3,\H);
\node at (0.9,3.22) {$x$};
\node at (0.9,2.32) {$g'(x)$};
\node at (0.9,0.9) {$g$};
\node at (2.25,3.22) {$-2$};
\node at (6.3,3.22) {$0$};
\node at (10.35,3.22) {$3$};
\node at (4.25,2.32) {$-$};\node at (6.3,2.32) {$0$};\node at (8.25,2.32) {$+$};
\node at (2.25,1.45) {$-3e^{-2}$};\node at (6.3,0.35) {$-1$};\node at (10.35,1.45) {$2e^3$};
\draw[->,line width=0.65pt] (2.75,1.35)--(5.8,0.45);
\draw[->,line width=0.65pt] (6.8,0.45)--(9.9,1.35);
\end{tikzpicture}
\end{center}}


\newcommand{\TableauSigneSujetUn}{%
\begin{center}
\begin{tikzpicture}[x=1cm,y=0.62cm,>=Latex]
\def\L{1.8}\def\W{10.8}\def\H{2.35}
\fill[black!8] (0,1.55) rectangle (\W,\H);
\fill[black!5] (0,0) rectangle (\L,\H);
\draw[tableline, line width=0.45pt] (0,0) rectangle (\W,\H);
\draw[tableline, line width=0.45pt] (\L,0)--(\L,\H);
\draw[tableline, line width=0.45pt] (0,1.55)--(\W,1.55);
\draw[tableline, line width=0.35pt] (4.8,0)--(4.8,\H);
\draw[tableline, line width=0.35pt] (7.8,0)--(7.8,\H);
\node at (0.9,1.95) {$x$};
\node at (0.9,0.75) {$f'(x)$};
\node at (2.1,1.95) {$-1$};
\node at (4.8,1.95) {$1$};
\node at (7.8,1.95) {$4$};
\node at (10.45,1.95) {$5$};
\node at (3.35,0.75) {$-$};
\node at (4.8,0.75) {$0$};
\node at (6.3,0.75) {$+$};
\node at (7.8,0.75) {$0$};
\node at (9.2,0.75) {$-$};
\end{tikzpicture}
\end{center}}

\newcommand{\TableauVariationSujetUn}{%
\begin{center}
\begin{tikzpicture}[x=1cm,y=0.62cm,>=Latex]
\def\L{1.8}\def\W{10.8}\def\H{3.65}
\fill[black!8] (0,2.82) rectangle (\W,\H);
\fill[black!5] (0,0) rectangle (\L,\H);
\draw[tableline, line width=0.45pt] (0,0) rectangle (\W,\H);
\draw[tableline, line width=0.45pt] (\L,0)--(\L,\H);
\draw[tableline, line width=0.45pt] (0,2.82)--(\W,2.82);
\draw[tableline, line width=0.45pt] (0,1.82)--(\W,1.82);
\draw[tableline, line width=0.35pt] (4.8,0)--(4.8,\H);
\draw[tableline, line width=0.35pt] (7.8,0)--(7.8,\H);
\node at (0.9,3.22) {$x$};
\node at (0.9,2.32) {$f'(x)$};
\node at (0.9,0.9) {$f$};
\node at (2.1,3.22) {$-1$};
\node at (4.8,3.22) {$1$};
\node at (7.8,3.22) {$4$};
\node at (10.45,3.22) {$5$};
\node at (3.35,2.32) {$-$};\node at (4.8,2.32) {$0$};\node at (6.3,2.32) {$+$};\node at (7.8,2.32) {$0$};\node at (9.2,2.32) {$-$};
\node at (4.8,0.35) {$f(1)$};\node at (7.8,1.45) {$f(4)$};
\draw[->,line width=0.65pt] (2.3,1.35)--(4.4,0.45);
\draw[->,line width=0.65pt] (5.2,0.45)--(7.4,1.35);
\draw[->,line width=0.65pt] (8.2,1.35)--(10.0,0.45);
\end{tikzpicture}
\end{center}}



% Tableaux TikZ spécifiques aux sujets v17
\newcommand{\TableauSigneGvDixSept}{%
\begin{center}
\begin{tikzpicture}[x=1cm,y=0.62cm,>=Latex]
\def\L{1.8}\def\W{12.2}\def\H{2.35}
\fill[black!8] (0,1.55) rectangle (\W,\H);
\fill[black!5] (0,0) rectangle (\L,\H);
\draw[tableline, line width=0.45pt] (0,0) rectangle (\W,\H);
\draw[tableline, line width=0.45pt] (\L,0)--(\L,\H);
\draw[tableline, line width=0.45pt] (0,1.55)--(\W,1.55);
\draw[tableline, line width=0.35pt] (5.0,0)--(5.0,\H);
\draw[tableline, line width=0.35pt] (8.2,0)--(8.2,\H);
\node at (0.9,1.95) {$x$};
\node at (0.9,0.75) {$g(x)$};
\node at (2.3,1.95) {$-\infty$};
\node at (5.0,1.95) {$1$};
\node at (8.2,1.95) {$4$};
\node at (11.75,1.95) {$+\infty$};
\node at (3.65,0.75) {$+$};
\node at (5.0,0.75) {$0$};
\node at (6.6,0.75) {$-$};
\node at (8.2,0.75) {$0$};
\node at (10.0,0.75) {$+$};
\end{tikzpicture}
\end{center}}

\newcommand{\TableauSigneFvDixSept}{%
\begin{center}
\begin{tikzpicture}[x=1cm,y=0.62cm,>=Latex]
\def\L{1.8}\def\W{10.8}\def\H{2.35}
\fill[black!8] (0,1.55) rectangle (\W,\H);
\fill[black!5] (0,0) rectangle (\L,\H);
\draw[tableline, line width=0.45pt] (0,0) rectangle (\W,\H);
\draw[tableline, line width=0.45pt] (\L,0)--(\L,\H);
\draw[tableline, line width=0.45pt] (0,1.55)--(\W,1.55);
\draw[tableline, line width=0.35pt] (6.0,0)--(6.0,\H);
\node at (0.9,1.95) {$x$};
\node at (0.9,0.75) {$f(x)$};
\node at (2.3,1.95) {$1$};
\node at (6.0,1.95) {$4$};
\node at (10.35,1.95) {$8$};
\node at (2.3,0.75) {$0$};
\node at (4.15,0.75) {$-$};
\node at (6.0,0.75) {$0$};
\node at (8.0,0.75) {$+$};
\end{tikzpicture}
\end{center}}

\newcommand{\TableauSigneDerivRationnelle}{%
\begin{center}
\begin{tikzpicture}[x=1cm,y=0.62cm,>=Latex]
\def\L{1.8}\def\W{10.8}\def\H{2.35}
\fill[black!8] (0,1.55) rectangle (\W,\H);
\fill[black!5] (0,0) rectangle (\L,\H);
\draw[tableline, line width=0.45pt] (0,0) rectangle (\W,\H);
\draw[tableline, line width=0.45pt] (\L,0)--(\L,\H);
\draw[tableline, line width=0.45pt] (0,1.55)--(\W,1.55);
\draw[tableline, line width=0.35pt] (5.8,0)--(5.8,\H);
\node at (0.9,1.95) {$x$};
\node at (0.9,0.75) {$f'(x)$};
\node at (2.3,1.95) {$1$};
\node at (5.8,1.95) {$2$};
\node at (10.35,1.95) {$8$};
\node at (4.05,0.75) {$-$};
\node at (5.8,0.75) {$0$};
\node at (8.0,0.75) {$+$};
\end{tikzpicture}
\end{center}}

\newcommand{\TableauVariationRationnelle}{%
\begin{center}
\begin{tikzpicture}[x=1cm,y=0.62cm,>=Latex]
\def\L{1.8}\def\W{10.8}\def\H{3.65}
\fill[black!8] (0,2.82) rectangle (\W,\H);
\fill[black!5] (0,0) rectangle (\L,\H);
\draw[tableline, line width=0.45pt] (0,0) rectangle (\W,\H);
\draw[tableline, line width=0.45pt] (\L,0)--(\L,\H);
\draw[tableline, line width=0.45pt] (0,2.82)--(\W,2.82);
\draw[tableline, line width=0.45pt] (0,1.82)--(\W,1.82);
\draw[tableline, line width=0.35pt] (5.8,0)--(5.8,\H);
\node at (0.9,3.22) {$x$};
\node at (0.9,2.32) {$f'(x)$};
\node at (0.9,0.9) {$f$};
\node at (2.3,3.22) {$1$};
\node at (5.8,3.22) {$2$};
\node at (10.35,3.22) {$8$};
\node at (4.05,2.32) {$-$};\node at (5.8,2.32) {$0$};\node at (8.0,2.32) {$+$};
\node at (2.3,1.45) {$0$};\node at (5.8,0.35) {$-1$};\node at (10.35,1.45) {$\frac72$};
\draw[->,line width=0.65pt] (2.8,1.35)--(5.35,0.45);
\draw[->,line width=0.65pt] (6.25,0.45)--(9.85,1.35);
\end{tikzpicture}
\end{center}}

\newcommand{\TableauVariationParamExpo}{%
\begin{center}
\begin{tikzpicture}[x=1cm,y=0.62cm,>=Latex]
\def\L{1.8}\def\W{12.4}\def\H{3.65}
\fill[black!8] (0,2.82) rectangle (\W,\H);
\fill[black!5] (0,0) rectangle (\L,\H);
\draw[tableline, line width=0.45pt] (0,0) rectangle (\W,\H);
\draw[tableline, line width=0.45pt] (\L,0)--(\L,\H);
\draw[tableline, line width=0.45pt] (0,2.82)--(\W,2.82);
\draw[tableline, line width=0.45pt] (0,1.82)--(\W,1.82);
\draw[tableline, line width=0.35pt] (6.2,0)--(6.2,\H);
\node at (0.9,3.22) {$x$};
\node at (0.9,2.32) {$f'(x)$};
\node at (0.9,0.9) {$f$};
\node at (2.25,3.22) {$-\infty$};
\node at (6.2,3.22) {$\frac37$};
\node at (12.0,3.22) {$+\infty$};
\node at (4.2,2.32) {$+$};\node at (6.2,2.32) {$0$};\node at (9.0,2.32) {$-$};
\node at (6.2,1.45) {$7e^{-3/7}$};
\draw[->,line width=0.65pt] (2.7,0.45)--(5.75,1.35);
\draw[->,line width=0.65pt] (6.65,1.35)--(11.45,0.45);
\end{tikzpicture}
\end{center}}

\newcommand{\TableauVariationConcentration}{%
\begin{center}
\begin{tikzpicture}[x=1cm,y=0.62cm,>=Latex]
\def\L{1.8}\def\W{10.8}\def\H{3.65}
\fill[black!8] (0,2.82) rectangle (\W,\H);
\fill[black!5] (0,0) rectangle (\L,\H);
\draw[tableline, line width=0.45pt] (0,0) rectangle (\W,\H);
\draw[tableline, line width=0.45pt] (\L,0)--(\L,\H);
\draw[tableline, line width=0.45pt] (0,2.82)--(\W,2.82);
\draw[tableline, line width=0.45pt] (0,1.82)--(\W,1.82);
\draw[tableline, line width=0.35pt] (5.8,0)--(5.8,\H);
\node at (0.9,3.22) {$t$};
\node at (0.9,2.32) {$C'(t)$};
\node at (0.9,0.9) {$C$};
\node at (2.3,3.22) {$0$};
\node at (5.8,3.22) {$1$};
\node at (10.35,3.22) {$5$};
\node at (4.05,2.32) {$+$};\node at (5.8,2.32) {$0$};\node at (8.0,2.32) {$-$};
\node at (2.3,0.35) {$0$};\node at (5.8,1.45) {$\frac4e$};\node at (10.35,0.35) {$20e^{-5}$};
\draw[->,line width=0.65pt] (2.8,0.45)--(5.35,1.35);
\draw[->,line width=0.65pt] (6.25,1.35)--(9.85,0.45);
\end{tikzpicture}
\end{center}}

\begin{document}
\begin{center}
{\LARGE\bfseries Corrigé -- Parcours de révision -- Dérivation}\\[1mm]
{\large Première générale -- spécialité mathématiques}\\[1mm]
{\normalsize Sans calculatrice}
\end{center}
\vspace{2mm}
\tableofcontents
\clearpage

\phantomsection
\addcontentsline{toc}{section}{Partie 1 -- Récapitulatif de cours}
\newgeometry{margin=6.5mm}
\pagestyle{empty}
\begin{landscape}
\pagecolor{fondpage}\color{encre}
\thispagestyle{empty}

\fontsize{7.85}{8.45}\selectfont

\titrepage

\noindent
\begin{minipage}[t]{0.242\linewidth}
\tagcol{bleu}{1. NOMBRE DÉRIVÉ}
\begin{panel}{bleu}

\begin{fichebox}{bleu}{Taux de variation}
Entre $a$ et $a+h$, le taux de variation de $f$ est :
\[
\frac{f(a+h)-f(a)}{h}.
\]
Il mesure une variation moyenne de la fonction.
\end{fichebox}

\begin{fichebox}{vert}{Nombre dérivé}
Si la limite existe, le nombre dérivé de $f$ en $a$ est :
\[
 f'(a)=\lim_{h\to0}\frac{f(a+h)-f(a)}{h}.
\]
Il donne la pente instantanée au point d'abscisse $a$.
\end{fichebox}

\begin{fichebox}{orange}{Lecture graphique}
$f'(a)$ est le coefficient directeur de la tangente à la courbe de $f$ au point d'abscisse $a$.

Si la tangente monte : $f'(a)>0$.

Si elle descend : $f'(a)<0$.

Si elle est horizontale : $f'(a)=0$.
\end{fichebox}

\begin{fichebox}{rouge}{Vocabulaire}
$f(a)$ est une ordonnée.

$f'(a)$ est un coefficient directeur.

Ne pas confondre la courbe et sa tangente.
\end{fichebox}

\end{panel}
\end{minipage}\hfill
\begin{minipage}[t]{0.242\linewidth}
\tagcol{vert}{2. TANGENTE}
\begin{panel}{vert}

\begin{fichebox}{vert}{Équation de la tangente}
Si $f$ est dérivable en $a$, l'équation de la tangente à la courbe de $f$ au point d'abscisse $a$ est :
\[
 y=f'(a)(x-a)+f(a).
\]
Ici, $f'(a)$ est le coefficient directeur et $A(a;f(a))$ est le point de contact.
\end{fichebox}

\begin{fichebox}{bleu}{Cas particuliers}
Tangente horizontale :
\[
 f'(a)=0.
\]
Tangente montante : coefficient directeur positif.

Tangente descendante : coefficient directeur négatif.
\end{fichebox}

\end{panel}
\end{minipage}\hfill
\begin{minipage}[t]{0.242\linewidth}
\tagcol{orange}{3. CALCULER UNE DÉRIVÉE}
\begin{panel}{orange}

\begin{fichebox}{orange}{Dérivées usuelles}
\[
\begin{array}{c|c}
f(x)&f'(x)\\ \hline
k&0\\
x&1\\
x^2&2x\\
x^3&3x^2\\
x^n&nx^{n-1}\\
e^x&e^x\\
\frac1x&-\frac1{x^2}\\
\sqrt{x}&\frac1{2\sqrt{x}}
\end{array}
\]
Sur les intervalles où les fonctions sont définies.
\end{fichebox}

\begin{fichebox}{bleu}{Opérations}
\[
(u+v)'=u'+v'
\]
\[
(ku)'=ku'
\]
Si $h(x)=g(ax+b)$, alors :
\[
h'(x)=a\,g'(ax+b).
\]
\[
(uv)'=u'v+uv'
\]
\[
\left(\frac{u}{v}\right)'=\frac{u'v-uv'}{v^2}
\quad (v\neq0).
\]
\end{fichebox}

\begin{fichebox}{rouge}{Erreurs fréquentes}
\[
(uv)'\neq u'v'.
\]
\[
\left(\frac{u}{v}\right)'\neq\frac{u'}{v'}.
\]
Toujours vérifier le domaine avant d'utiliser une formule.
\end{fichebox}

\end{panel}
\end{minipage}\hfill
\begin{minipage}[t]{0.242\linewidth}
\tagcol{violet}{4. VARIATIONS}
\begin{panel}{violet}

\begin{fichebox}{violet}{Lien signe / variations}
Sur un intervalle :
\[
\begin{array}{c|c}
f'(x)>0&f\text{ croissante}\\
f'(x)<0&f\text{ décroissante}\\
f'(x)=0&\text{tangente horizontale possible}
\end{array}
\]
\end{fichebox}

\begin{fichebox}{vert}{Méthode : tableau de variations}
\begin{enumerate}[leftmargin=*,itemsep=0.1em]
\item Calculer $f'(x)$.
\item Factoriser si possible.
\item Étudier le signe de $f'(x)$.
\item En déduire les variations de $f$.
\end{enumerate}
\end{fichebox}

\begin{fichebox}{orange}{Extremum}
Si $f'$ change de signe de $+$ vers $-$, alors $f$ admet un maximum local.

Si $f'$ change de signe de $-$ vers $+$, alors $f$ admet un minimum local.
\end{fichebox}

\end{panel}
\end{minipage}

\end{landscape}
\nopagecolor
\restoregeometry
\pagestyle{fancy}
\clearpage
\section*{Partie 1 -- Corrigé des questions flash}\addcontentsline{toc}{section}{Partie 1 -- Corrigé des questions flash}
\begin{blocreponses}{Réponses -- Questions flash}
  \item Il représente le coefficient directeur de la tangente à la courbe de $f$ au point d'abscisse $a$.
  \item $\dfrac{f(a+h)-f(a)}{h}$.
  \item Elle est montante.
  \item Elle est descendante.
  \item Elle est horizontale.
  \item $0$.
  \item $f'(x)=1$.
  \item $f'(x)=2x$.
  \item $f'(x)=3x^2$.
  \item $f'(x)=nx^{n-1}$ pour $n$ entier naturel non nul.
  \item $f'(x)=-\dfrac1{x^2}$.
  \item $f'(x)=\dfrac1{2\sqrt{x}}$ sur $]0;+\infty[$.
  \item $(u+v)'=u'+v'$.
  \item $(ku)'=ku'$.
  \item $(uv)'=u'v+uv'$.
  \item $(u/v)'=\dfrac{u'v-uv'}{v^2}$ lorsque $v\neq0$.
  \item $f$ est croissante sur cet intervalle.
  \item $f$ est décroissante sur cet intervalle.
  \item On calcule $f'$, on étudie son signe, puis on dresse le tableau de variations de $f$.
  \item Un maximum local.
  \item Un minimum local.
  \item $y=f'(a)(x-a)+f(a)$.
  \item $f'(a)$ est le coefficient directeur de la tangente.
  \item $y=3(x-2)+5$, c'est-à-dire $y=3x-1$.
  \item $h'(x)=2g'(2x-1)$.
  \item $f'(x)=6x-5$.
  \item $f'(x)=-6x^2+4$.
  \item On peut développer : $f(x)=x^2-x-2$, donc $f'(x)=2x-1$.
  \item On peut développer : $f(x)=3x^3+x^2$, donc $f'(x)=9x^2+2x$.
  \item Croire que $(uv)'=u'v'$, ce qui est faux.
  \item Croire que $(u/v)'=u'/v'$, ce qui est faux.
  \item $f'(x)=e^x$.
  \item $f'(x)=2e^{2x-1}$.
  \item $f'(x)=3e^x-2$.
  \item $f'(x)>0$.
  \item $f'(x)<0$.
\end{blocreponses}
\clearpage
\section*{Partie 2 -- Corrigé du QCM}\addcontentsline{toc}{section}{Partie 2 -- Corrigé du QCM}
\begin{correctionbox}{Réponses du QCM}
\begin{center}\begin{tabular}{c|cccccccc} Question & 1 & 2 & 3 & 4 & 5 & 6 & 7 & 8 \\ \hline Réponse & b & a & c & c & b & b & a & b\end{tabular}\end{center}
\end{correctionbox}
\clearpage


\section*{Partie 3 -- Corrigés des sujets type bac / E3C}
\addcontentsline{toc}{section}{Partie 3 -- Corrigés des sujets type bac / E3C}

\titreSujet{Sujet 1 -- Fonction rationnelle, signe et variations}{inspiré du sujet zéro officiel : signe d'un trinôme, fonction rationnelle, dérivée et variations}
\begin{correctionbox}{1. Factorisation et signe de $g$}
On a :
\[
 g(x)=x^2-5x+4.
\]
On cherche deux nombres dont la somme vaut $5$ et le produit vaut $4$ : ce sont $1$ et $4$. Donc :
\[
 \boxed{g(x)=(x-1)(x-4)}.
\]
Le coefficient dominant est positif, donc $g$ est positive à l'extérieur des racines et négative entre les racines.

\TableauSigneGvDixSept
\end{correctionbox}

\begin{correctionbox}{2. Lien entre $f$ et $g$}
Pour tout $x\in[1;8]$, on a $x>0$. On calcule :
\[
 \frac{g(x)}x=\frac{x^2-5x+4}{x}.
\]
Donc :
\[
 \frac{g(x)}x=x-5+\frac4x.
\]
Ainsi :
\[
 \boxed{f(x)=\frac{g(x)}x}.
\]
\end{correctionbox}

\begin{correctionbox}{3. Signe de $f$ sur $[1;8]$}
Sur $[1;8]$, on a $x>0$. Donc $f(x)=\dfrac{g(x)}x$ a le même signe que $g(x)$.

D'après le tableau de signe de $g$ :
\[
 f(1)=0,\qquad f(4)=0,
\]
$f(x)<0$ sur $]1;4[$ et $f(x)>0$ sur $]4;8]$.

\TableauSigneFvDixSept
\end{correctionbox}

\begin{correctionbox}{4. Calcul de $f'(x)$}
On écrit :
\[
 f(x)=x-5+\frac4x.
\]
Pour $x\in[1;8]$, on obtient :
\[
 f'(x)=1+0-\frac4{x^2}.
\]
Donc :
\[
 \boxed{f'(x)=1-\frac4{x^2}}.
\]
\end{correctionbox}

\begin{correctionbox}{5. Forme factorisée de $f'(x)$}
On part de :
\[
 f'(x)=1-\frac4{x^2}.
\]
On met au même dénominateur :
\[
 f'(x)=\frac{x^2-4}{x^2}.
\]
Puis on factorise le numérateur :
\[
 x^2-4=(x-2)(x+2).
\]
Ainsi :
\[
 \boxed{f'(x)=\frac{(x-2)(x+2)}{x^2}}.
\]
\end{correctionbox}

\begin{correctionbox}{6. Signe de $f'$}
Sur $[1;8]$, on a $x^2>0$ et $x+2>0$. Le signe de $f'(x)$ est donc celui de $x-2$.

\TableauSigneDerivRationnelle
\end{correctionbox}

\begin{correctionbox}{7. Tableau de variations de $f$}
On calcule les valeurs utiles :
\[
 f(1)=1-5+4=0,
\]
\[
 f(2)=2-5+2=-1,
\]
\[
 f(8)=8-5+\frac48=3+\frac12=\frac72.
\]
On obtient :

\TableauVariationRationnelle
\end{correctionbox}

\begin{correctionbox}{8. Minimum de $f$}
D'après le tableau de variations, $f$ atteint son minimum sur $[1;8]$ en $x=2$.

Ce minimum vaut :
\[
 \boxed{-1}.
\]
\end{correctionbox}

\begin{correctionbox}{9. Interprétation graphique}
Le signe de $f(x)$ donne la position de la courbe de $f$ par rapport à l'axe des abscisses.

Ainsi, la courbe coupe l'axe des abscisses en $x=1$ et $x=4$, elle est sous l'axe sur $]1;4[$, puis au-dessus de l'axe sur $]4;8]$.
\end{correctionbox}

\clearpage
\titreSujet{Sujet 2 -- Fonction avec paramètre et tangente}{inspiré d'un sujet E3C : tangente, paramètres, produit avec exponentielle et variations}
\begin{correctionbox}{1. Détermination de $b$}
La courbe passe par $A(0;4)$, donc :
\[
 f(0)=4.
\]
Or :
\[
 f(0)=(a\times0+b)e^0=b.
\]
Donc :
\[
 \boxed{b=4}.
\]
\end{correctionbox}

\begin{correctionbox}{2. Coefficient directeur de la tangente}
La tangente en $A$ passe par $A(0;4)$ et $B(1;7)$. Son coefficient directeur est :
\[
 m=\frac{7-4}{1-0}=3.
\]
\end{correctionbox}

\begin{correctionbox}{3. Valeur de $f'(0)$}
Le nombre $f'(0)$ est le coefficient directeur de la tangente à la courbe au point d'abscisse $0$.

Donc :
\[
 \boxed{f'(0)=3}.
\]
\end{correctionbox}

\begin{correctionbox}{4. Calcul de $f'(x)$}
On pose :
\[
 u(x)=ax+b \qquad \text{et} \qquad v(x)=e^{-x}.
\]
Alors :
\[
 u'(x)=a \qquad \text{et} \qquad v'(x)=-e^{-x}.
\]
D'après la formule de dérivée d'un produit :
\[
 f'(x)=u'(x)v(x)+u(x)v'(x).
\]
Donc :
\[
 f'(x)=ae^{-x}-(ax+b)e^{-x}.
\]
Ainsi :
\[
 \boxed{f'(x)=(a-ax-b)e^{-x}}.
\]
\end{correctionbox}

\begin{correctionbox}{5. Détermination de $a$}
Comme $b=4$, on a :
\[
 f'(x)=(a-ax-4)e^{-x}.
\]
Donc :
\[
 f'(0)=(a-4)e^0=a-4.
\]
Or $f'(0)=3$. Ainsi :
\[
 a-4=3,
\]
soit :
\[
 \boxed{a=7}.
\]
Donc :
\[
 \boxed{f(x)=(7x+4)e^{-x}}.
\]
\end{correctionbox}

\begin{correctionbox}{6. Signe de $f'$}
Avec $a=7$ et $b=4$, on obtient :
\[
 f'(x)=(7-7x-4)e^{-x}=(3-7x)e^{-x}.
\]
Comme $e^{-x}>0$ pour tout réel $x$, le signe de $f'(x)$ est celui de $3-7x$.

On a :
\[
 3-7x=0 \quad \Longleftrightarrow \quad x=\frac37.
\]
Donc $f'(x)>0$ si $x<\frac37$ et $f'(x)<0$ si $x>\frac37$.
\end{correctionbox}

\begin{correctionbox}{7. Tableau de variations de $f$}
On calcule la valeur au point critique :
\[
 f\left(\frac37\right)=\left(7\times\frac37+4\right)e^{-3/7}.
\]
Donc :
\[
 f\left(\frac37\right)=7e^{-3/7}.
\]
On obtient le tableau de variations suivant :

\TableauVariationParamExpo
\end{correctionbox}

\begin{correctionbox}{8. Maximum de $f$}
D'après le tableau de variations, $f$ admet un maximum sur $\R$ pour :
\[
 x=\frac37.
\]
Ce maximum vaut :
\[
 \boxed{7e^{-3/7}}.
\]
\end{correctionbox}

\begin{correctionbox}{9. Tangente au point $A$}
On connaît $A(0;4)$ et le coefficient directeur de la tangente : $m=3$.

Avec la formule de tangente :
\[
 y=f'(0)(x-0)+f(0),
\]
on obtient :
\[
 y=3x+4.
\]
Ainsi :
\[
 \boxed{y=3x+4}.
\]
\end{correctionbox}

\clearpage
\titreSujet{Sujet 3 -- Concentration et seuil réglementaire}{inspiré de sujets E3C : dérivée en contexte, maximum et interprétation}
\begin{correctionbox}{1. Valeur initiale}
\[
 C(0)=4\times0\times e^0=0.
\]
Donc :
\[
 \boxed{C(0)=0}.
\]
\end{correctionbox}

\begin{correctionbox}{2. Calcul de $C'(t)$}
On écrit :
\[
 C(t)=4t e^{-t}.
\]
On pose $u(t)=4t$ et $v(t)=e^{-t}$. Alors :
\[
 u'(t)=4 \qquad \text{et} \qquad v'(t)=-e^{-t}.
\]
D'après la formule de dérivée d'un produit :
\[
 C'(t)=u'(t)v(t)+u(t)v'(t).
\]
Donc :
\[
 C'(t)=4e^{-t}-4t e^{-t}.
\]
Ainsi :
\[
 \boxed{C'(t)=4(1-t)e^{-t}}.
\]
\end{correctionbox}

\begin{correctionbox}{3. Signe de $C'$}
Pour tout $t\in[0;5]$, on a $e^{-t}>0$ et $4>0$.

Le signe de $C'(t)=4(1-t)e^{-t}$ est donc celui de $1-t$.

Ainsi, $C'(t)>0$ sur $[0;1[$, $C'(1)=0$, puis $C'(t)<0$ sur $]1;5]$.
\end{correctionbox}

\begin{correctionbox}{4. Tableau de variations de $C$}
On calcule :
\[
 C(0)=0,
\]
\[
 C(1)=4e^{-1}=\frac4e,
\]
\[
 C(5)=20e^{-5}.
\]
On obtient le tableau suivant :

\TableauVariationConcentration
\end{correctionbox}

\begin{correctionbox}{5. Instant du maximum}
D'après le tableau de variations, la concentration est maximale pour :
\[
 \boxed{t=1}.
\]
\end{correctionbox}

\begin{correctionbox}{6. Valeur exacte de la concentration maximale}
La concentration maximale vaut :
\[
 \boxed{\frac4e} \quad \text{mg/L}.
\]
\end{correctionbox}

\begin{correctionbox}{7. Comparaison avec le seuil de $1$ mg/L}
On admet que :
\[
 2<e<3.
\]
Alors $e<3<4$, donc :
\[
 \frac4e>1.
\]
La concentration maximale est donc supérieure à $1$ mg/L.
\end{correctionbox}

\begin{correctionbox}{8. Valeur au bout de $3$ h}
On a :
\[
 C(3)=4\times3\times e^{-3}=12e^{-3}=\frac{12}{e^3}.
\]
On admet que $e^3>12$. Donc :
\[
 \frac{12}{e^3}<1.
\]
Ainsi :
\[
 \boxed{C(3)<1}.
\]
Au bout de $3$ h, la réglementation est de nouveau respectée.
\end{correctionbox}

\begin{correctionbox}{9. Interprétation}
Au départ, la concentration est nulle. Elle augmente jusqu'à atteindre son maximum au bout d'une heure, puis elle diminue.

Comme le maximum est supérieur à $1$ mg/L, le seuil réglementaire est dépassé à un certain moment. Comme $C(3)<1$, ce seuil n'est plus dépassé au bout de $3$ h.
\end{correctionbox}

\end{document}
